\section{Boundary metrics and their local profiles}\label{sec:frames}
The representation in this section is the frame construction of
M\"uller--P\'or--Sereni \cite[Lemmas 7--8 and Theorem 13]{muller}.
We include its short proof to make the completeness of the finite
certificates explicit. The new contribution here is the harmonic-boundary
exclusion obtained from this representation.

\begin{lemma}[Boundary profiles]\label{lem:profiles}
Let $B=\{b_1,\ldots,b_k\}$ contain the CEJZ boundary of a connected graph.
Write $M_{ij}=d(b_i,b_j)$ and $f(x)=(d(x,b_1),\ldots,d(x,b_k))$. Then
\begin{align*}
 d(x,y)&=\|f(x)-f(y)\|_\infty,\\
 f_i(x)+f_j(x)&\ge M_{ij},\\
 \text{for each }i\text{ there is }j\text{ with }f_i(x)+f_j(x)&=M_{ij}.
\end{align*}
In particular the profiles are injective, and adjacency is forced:
$x\sim y$ if and only if their profile distance is one.
Moreover $\stab(u)=\max_{b\in B}\stab(u,b)$.
\end{lemma}
\begin{proof}
Extend a shortest $x$--$y$ path by steps that increase the distance from
$x$, until no such step exists. Its final vertex $b$ is CEJZ boundary
with witness $x$, and $d(x,b)=d(x,y)+d(y,b)$. The triangle inequality
gives the reverse bound for every coordinate, proving the norm identity.
The sum inequalities are triangle inequalities. Extending a shortest
$b_i$--$x$ path in the same way gives the required tight partner $b_j$.
For the last statement extend a shortest $u$--$t$ path to $b\in B$.
For every $w\sim u$,
\[
 d(w,b)\le d(w,t)+d(t,b),\qquad
 d(u,b)=d(u,t)+d(t,b),
\]
so $\stab(u,b)\ge\stab(u,t)$. Maximizing proves the identity.
\end{proof}
The difference inequalities $|f_i-f_j|\le M_{ij}$ also follow from the
displayed sum and tightness conditions: choose $\ell$ tight with $i$,
and use $f_i+f_\ell=M_{i\ell}\le M_{ij}+M_{j\ell}
\le M_{ij}+f_j+f_\ell$. Thus the two independent implementations may
include these inequalities explicitly or derive them implicitly.

\begin{lemma}[No long pendant metric segment]\label{lem:pendantmetric}
In a two-connected graph, let $B$ contain its CEJZ boundary and have at
least three vertices. For every $b\in B$,
\[
 \min_{j\ne k,\ b_j,b_k\ne b}
 \bigl(d(b,b_j)+d(b,b_k)-d(b_j,b_k)\bigr)\le1.
\]
Equivalently, the usual pendant length of each designated boundary point
in its finite metric is zero or one half.
\end{lemma}
\begin{proof}
If that minimum were at least two, let $x$ be any neighbor of $b$.
Its coordinate at $b$ equals one, and $f_j(x)\ge d(b,b_j)-1$ for all
other coordinates. If $f_j(x)\ge d(b,b_j)$, no other coordinate $k$
can be tight with $j$: for $k\ne b,j$ the assumed inequality gives
$f_j(x)+f_k(x)>d(b_j,b_k)$, for $k=b$ this is immediate, and
$k=j$ cannot be tight because $f_j(x)>0$.
Thus every $j\ne b$ has $f_j(x)=d(b,b_j)-1$.
Every neighbor of $b$ has the identical profile. Injectivity of profiles
forces degree one, contradicting two-connectivity.
\end{proof}
This restriction is independent of the diameter. It bounds the five
pendant coefficients used in the uniform finite reduction below.
